\section{Introduction}
\label{sec:introduction}

\IEEEPARstart{C}{ontinuous}, non-destructive observation of adherent cell
cultures inside perfused bioreactor chips---for example hFOB~1.19 osteoblasts
and hMEC~1 endothelial cells under pressure and flow test series---is a
recurring need in tissue-engineering and mechanobiology experiments. Today it is
served either by \emph{manual microscopy}, which removes the chip from the
incubator, is operator-dependent and offers low temporal resolution, or by
\emph{commercial live-cell imagers}, which are expensive, closed and typically
observe a single vessel at a time.

This thesis designs, builds and validates a \emph{mesoscope}: an
OpenFlexure-derived microscope running on a Raspberry~Pi~5 with a Hailo-8
neural-processing unit (NPU) and a Pi High-Quality (HQ) camera. The microscope
is laid horizontally and mounted on a linear rail so that a single imaging head
scans a rack of bioreactor chips in series, and it evaluates cell-level images
on-device without external compute infrastructure.

\subsection{Problem statement}
The core problem is to observe several perfused bioreactor chips repeatedly,
inside the incubator, through their small optical windows, at a spatial and
temporal resolution adequate for automated cell analysis, at a cost far below
commercial live-cell imaging systems.

\subsection{Research questions}
\begin{itemize}
\item \textbf{RQ1 (System).} Can an open-source, sub-\num{500}~EUR
  microscope built from OpenFlexure, a Pi HQ camera and 3D-printed parts deliver
  image quality (resolution, field of view, focus stability) sufficient for
  automated cell analysis in a bioreactor chip through its observation window?
\item \textbf{RQ2 (Automation).} Can a single imaging head on a linear X-rail
  revisit $N$ bioreactors autonomously with positioning repeatability and
  autofocus reliability adequate for time-lapse monitoring over days inside an
  incubator?
\item \textbf{RQ3 (Edge AI).} Can a cell-detection/segmentation model trained on
  data annotated in CVAT be deployed on the Hailo-8 NPU with accuracy
  (mAP / IoU / count error) and throughput sufficient for on-device,
  near-real-time analysis---and how does it compare to inference on the Pi CPU
  and to manual counting?
\end{itemize}

\subsection{Hypotheses and success criteria}
\label{sec:hypotheses}
The following criteria are the current targets; they are finalised in
milestone~M1 (see Section~\ref{sec:evaluation}).
\begin{itemize}
\item \textbf{H1.} Measured lateral resolution $\leq \SI{1.5}{\micro\meter}$
  (USAF~1951 target; Rayleigh limit $\approx \SI{1.1}{\micro\meter}$ at
  NA~$0.30$) and field of view $\geq \SI{0.5}{\milli\meter\squared}$ with the
  ZEISS Plan-Neofluar $10\times/0.30$ objective and a \SI{125}{\milli\meter}
  tube lens.
\item \textbf{H2.} Rail positioning repeatability
  $\leq \pm\SI{10}{\micro\meter}$; autofocus success rate $\geq \SI{95}{\percent}$
  across $\geq 3$ chips over \SI{72}{\hour}.
\item \textbf{H3.} Cell-detection mAP@0.5 $\geq 0.8$ on held-out chips; Hailo-8
  inference $\geq 10\times$ faster than the Pi CPU at equal accuracy
  ($\pm\SI{2}{\percent}$ mAP).
\end{itemize}

\subsection{Contributions}
\label{sec:contributions}
\begin{enumerate}
\item A reproducible port of the OpenFlexure Server~v3 to Raspberry~Pi~5 /
  Trixie via Ansible, including documented fixes for the Pi~5 image-signal
  processor (ISP), camera stack and Hailo driver.
\item A C-mount HQ-camera optics module with a \SI{125}{\milli\meter} (and
  \SI{150}{\milli\meter}) tube lens for the IMX477 sensor, with full optical
  derivation and printable CAD.
\item A tilted, rail-mounted imaging head with a Z-focus-only body and a
  parametric bioreactor rack---a new form factor for serial monitoring of
  multiple chips.
\item An all-Python GPIO stage driver (focus stepper + rail stepper) replacing
  the Sangaboard microcontroller. \todo{confirm motor architecture, milestone
  M2.}
\item A CVAT-annotated cell dataset from bioreactor-chip images and a
  Hailo-8-deployed detection model with a CPU-versus-NPU evaluation.
\item Quantitative validation of the complete system against the stated
  requirements.
\end{enumerate}

\subsection{Thesis structure}
Section~\ref{sec:background} surveys live-cell imaging, open-hardware microscopy
and edge AI. Section~\ref{sec:requirements} derives requirements and the system
concept. Sections~\ref{sec:optical}--\ref{sec:software} present the optical,
mechanical, electronic and software design. Section~\ref{sec:edgeai} describes
the edge-AI pipeline. Section~\ref{sec:evaluation} evaluates the system,
Section~\ref{sec:discussion} answers RQ1--RQ3, and
Section~\ref{sec:conclusion} concludes.
